\newglossaryentry{overfitting}
{
    name={overfitting},
    description={When a model learns the training data too well, including noise, and performs poorly on new data}
}

\newglossaryentry{underfitting}
{
    name={underfitting},
    description={When a model is too simple and cannot capture the underlying patterns in the data}
}

\newglossaryentry{accuracy}
{
    name={accuracy},
    description={Percentage of correct predictions out of all predictions}
}

\newglossaryentry{generalisation}
{
    name={generalisation},
    description={Ability of a model to perform well on unseen data}
}

\newglossaryentry{feature}
{
    name={feature},
    description={An individual variable from the data}
}

\newglossaryentry{label}
{
    name={label},
    description={The output variable associated with a data point. In this case, which team won}
}

\newglossaryentry{intent}{
    name={intent},
    description={An intent represents a function in the code. It can have multiple utterances attached to it to trigger it.}
}

\newglossaryentry{utterance}{
    name={utterance},
    description={A predefined phrase that can be spoken to trigger the attached intent.}
}

\newglossaryentry{hyper-parameter}{
    name={hyper-parameter},
    description={A variable used for configuration of the model, set before training takes place.}
}

\newglossaryentry{classification}{
    name={classification},
    description={A type of machine learning problem where the predictions fall into discrete classes such as "Dog" or "Cat".}
}

\newglossaryentry{regression}{
    name={regression},
    description={A type of machine learning problem where the prediction is a continuous value such as expected salary.}
}

\newglossaryentry{ensemble}{
    name={ensemble},
    description={A model that combines multiple individual machine learning models to produce more accurate and robust predictions than a single model could alone.}
}

\newglossaryentry{regularise}{
    name={regularise},
    description={To employ a penalty to the model if it becomes too complex. Useful when models are prone to overfitting; the 'cost' hyper-parameter can be tuned to ensure that the model generalises well.}
}

% Synonyms using "see" references
\newglossaryentry{regularised}{
    name={regularised},
    description={see \glsentryname{regularise}},
    sort=regularise
}

\newglossaryentry{regularisation}{
    name={regularisation},
    description={see \glsentryname{regularise}},
    sort=regularise
}

\newacronym{ml}{ML}{Machine Learning}
\newacronym{vct}{VCT}{Valorant Champions Tour}
\newacronym{aws}{AWS}{Amazon Web Services}
\newacronym{svm}{SVM}{Support Vector Machine}
\newacronym{knn}{KNN}{K-Nearest Neighbours}
\newacronym{gbt}{GBT}{Gradient Boosted Trees}
\newacronym{vui}{VUI}{Voice User Interface}
\newacronym{sdk}{SDK}{Software Development Kit}
\newacronym{avs}{AVS}{Alexa Voice Services}
\newacronym{pii}{PII}{Personal Identifiable Information}